% GENERATED FILE. Edit canonical lessons, then run npm run generate:books.
% canonical-source-hash: fnv1a64:b0e8e54195c36426
% canonical-lessons: JA-C319-honne, JA-C319-tatemae, JA-C319-mitomeru, JA-C319-gokai-suru, JA-C319-honki

\chapter{What People Really Think}
\label{ch:ja-what-people-really-think}
\hlchaptermodality{319}

\begin{chapteropening}
\textbf{By the end of this chapter:} \emph{I can tell what someone says from what they really think.}

Complete the last lesson of chapter 319: I can tell what someone says from what they really think.
\end{chapteropening}

\section[\texorpdfstring{\ja{ほんね} (honne)}{ほんね (honne)}]{\ja{ほんね} (honne) — one's real feelings; what one really thinks}

\label{lesson:JA-C319-honne}



\begin{quote}
Before the new one: say the Japanese for a meaning, then the Japanese for for what reason.
\end{quote}

\subsection*{You'll want to know: \ja{ほんね}}
\textbf{\ja{ほんね}} — \emph{honne} — \textquotedblleft{}one's real feelings; what one really thinks\textquotedblright{}.

\begin{quote}
\textbf{\ja{しんゆうには} \ja{ほんねを} \ja{はなします。}} — \emph{shin'yū ni wa honne o hanashimasu} — \textquotedblleft{}I tell my best friend what I really think.\textquotedblright{}
\end{quote}

\begin{cousinweb}
\textbf{\ja{ほんね}} (\emph{honne}) is built from a Sino-Japanese part and a native part: \emph{hon}, \textquotedblleft{}true\textquotedblright{}, an on'yomi borrowed from Chinese (Middle Chinese \emph{pwonX}), and \emph{ne}, \textquotedblleft{}sound\textquotedblright{}, a native kun'yomi (KANJIDIC2; JMdict; Middle Chinese in Baxter's transcription).
\end{cousinweb}

\begin{grammarlens}[title={What you've built}]
One more word for giving reasons and opinions.
\end{grammarlens}

\subsection*{Your turn}
\begin{itemize}
\raggedright
  \item \emph{Say it:} \emph{honne}
  \item \emph{Say it:} \emph{honne}, once more
  \item \emph{Say it:} \emph{shin'yū ni wa honne o hanashimasu} — I tell my best friend what I really think.
  \item \emph{Recall:} say the Japanese for a meaning, then the Japanese for for what reason, then say \emph{honne} again
\end{itemize}

\begin{tcolorbox}[breakable,colback=teal!4,colframe=teal!35!black,title={Before you move on}]
What does \ja{ほんね} mean? (\textquotedblleft{}One's real feelings; what one really thinks\textquotedblright{}.) Say it once more.
\end{tcolorbox}
\section[\texorpdfstring{\ja{たてまえ} (tatemae)}{たてまえ (tatemae)}]{\ja{たてまえ} (tatemae) — the official line; what one is expected to say}

\label{lesson:JA-C319-tatemae}



\begin{quote}
Before the new one: say the Japanese for for what reason, then the Japanese for one's real feelings.
\end{quote}

\subsection*{You'll want to know: \ja{たてまえ}}
\textbf{\ja{たてまえ}} — \emph{tatemae} — \textquotedblleft{}the official line; what one is expected to say\textquotedblright{}.

\begin{quote}
\textbf{\ja{じょうしが} \ja{いたから、たてまえを} \ja{言}\ja{いました。}} — \emph{jōshi ga ita kara, tatemae o iimashita} — \textquotedblleft{}My boss was there, so I said what I was expected to say.\textquotedblright{}
\end{quote}

\begin{cousinweb}
\textbf{\ja{たてまえ}} (\emph{tatemae}) is a native Japanese word (wago), inherited rather than borrowed from Chinese: it is built from \emph{tate}, \textquotedblleft{}building (tateru)\textquotedblright{} and \emph{mae}, \textquotedblleft{}front\textquotedblright{}, and its kanji were attached later as kun'yomi (JMdict).
\end{cousinweb}

\begin{grammarlens}[title={What you've built}]
One more word for giving reasons and opinions.
\end{grammarlens}

\subsection*{Your turn}
\begin{itemize}
\raggedright
  \item \emph{Say it:} \emph{tatemae}
  \item \emph{Say it:} \emph{tatemae}, once more
  \item \emph{Say it:} \emph{jōshi ga ita kara, tatemae o iimashita} — My boss was there, so I said what I was expected to say.
  \item \emph{Recall:} say the Japanese for for what reason, then the Japanese for one's real feelings, then say \emph{tatemae} again
\end{itemize}

\begin{tcolorbox}[breakable,colback=teal!4,colframe=teal!35!black,title={Before you move on}]
What does \ja{たてまえ} mean? (\textquotedblleft{}The official line; what one is expected to say\textquotedblright{}.) Say it once more.
\end{tcolorbox}
\section[\texorpdfstring{\ja{みとめる} (mitomeru)}{みとめる (mitomeru)}]{\ja{みとめる} (mitomeru) — to admit; to recognize}

\label{lesson:JA-C319-mitomeru}



\begin{quote}
Before the new one: say the Japanese for one's real feelings, then the Japanese for the official line.
\end{quote}

\subsection*{You'll want to know: \ja{みとめる}}
\textbf{\ja{みとめる}} — \emph{mitomeru} — \textquotedblleft{}to admit; to recognize\textquotedblright{}.

\begin{quote}
\textbf{\ja{あには} \ja{まちがえたので、それを} \ja{みとめました。}} — \emph{ani wa machigaeta node, sore o mitomemashita} — \textquotedblleft{}My brother had made a mistake, so he admitted it.\textquotedblright{}
\end{quote}

Polite \textbf{\ja{みとめます}}, \emph{mitomemasu}; an ichidan verb. Past (polite): \textbf{\ja{あには} \ja{みとめました。}}, \emph{ani wa mitomemashita}.

\begin{cousinweb}
\textbf{\ja{みとめる}} (\emph{mitomeru}) is a native Japanese word (wago), inherited rather than borrowed from Chinese, and its older origin is not settled. Its kanji was attached to it later as a kun'yomi, a native reading given to a Chinese character (JMdict, KANJIDIC2).
\end{cousinweb}

\begin{grammarlens}[title={What you've built}]
One more word for giving reasons and opinions.
\end{grammarlens}

\subsection*{Your turn}
\begin{itemize}
\raggedright
  \item \emph{Say it:} \emph{mitomeru}
  \item \emph{Say it:} \emph{mitomeru}, once more
  \item \emph{Say it:} \emph{ani wa machigaeta node, sore o mitomemashita} — My brother had made a mistake, so he admitted it.
  \item \emph{Recall:} say the Japanese for one's real feelings, then the Japanese for the official line, then say \emph{mitomeru} again
\end{itemize}

\begin{tcolorbox}[breakable,colback=teal!4,colframe=teal!35!black,title={Before you move on}]
What does \ja{みとめる} mean? (\textquotedblleft{}To admit; to recognize\textquotedblright{}.) Say it once more.
\end{tcolorbox}
\section[\texorpdfstring{\ja{ごかいする} (gokai suru)}{ごかいする (gokai suru)}]{\ja{ごかいする} (gokai suru) — to misunderstand}

\label{lesson:JA-C319-gokai-suru}



\begin{quote}
Before the new one: say the Japanese for the official line, then the Japanese for to admit.
\end{quote}

\subsection*{You'll want to know: \ja{ごかいする}}
\textbf{\ja{ごかいする}} — \emph{gokai suru} — \textquotedblleft{}to misunderstand\textquotedblright{}.

\begin{quote}
\textbf{\ja{せつめいを} \ja{まちがえたので、みんな} \ja{ごかいしました。}} — \emph{setsumei o machigaeta node, minna gokai shimashita} — \textquotedblleft{}I got the explanation wrong, so everyone misunderstood.\textquotedblright{}
\end{quote}

Polite \textbf{\ja{ごかいします}}, \emph{gokaishimasu}; a \ja{する} verb. Past (polite): \textbf{\ja{ともだちが} \ja{ごかいしました。}}, \emph{tomodachi ga gokai shimashita}.

\begin{cousinweb}
\textbf{\ja{ごかいする}} (\emph{gokai suru}) is built from a Sino-Japanese noun and the native verb \emph{suru}, \textquotedblleft{}to do\textquotedblright{}. The noun, \emph{gokai}, is two Chinese characters read with their on'yomi, readings borrowed from Chinese: \emph{go}, \textquotedblleft{}mistake\textquotedblright{} (Middle Chinese \emph{nguH}) and \emph{kai}, \textquotedblleft{}understanding\textquotedblright{} (\emph{keaX}) (KANJIDIC2; JMdict; Middle Chinese in Baxter's transcription).
\end{cousinweb}

\begin{grammarlens}[title={What you've built}]
One more word for giving reasons and opinions.
\end{grammarlens}

\subsection*{Your turn}
\begin{itemize}
\raggedright
  \item \emph{Say it:} \emph{gokai suru}
  \item \emph{Say it:} \emph{gokai suru}, once more
  \item \emph{Say it:} \emph{setsumei o machigaeta node, minna gokai shimashita} — I got the explanation wrong, so everyone misunderstood.
  \item \emph{Recall:} say the Japanese for the official line, then the Japanese for to admit, then say \emph{gokai suru} again
\end{itemize}

\begin{tcolorbox}[breakable,colback=teal!4,colframe=teal!35!black,title={Before you move on}]
What does \ja{ごかいする} mean? (\textquotedblleft{}To misunderstand\textquotedblright{}.) Say it once more.
\end{tcolorbox}
\section[\texorpdfstring{\ja{ほんき} (honki)}{ほんき (honki)}]{\ja{ほんき} (honki) — serious; in earnest}

\label{lesson:JA-C319-honki}



\begin{quote}
Before the new one: say the Japanese for to admit, then the Japanese for to misunderstand.
\end{quote}

\subsection*{You'll want to know: \ja{ほんき}}
\textbf{\ja{ほんき}} — \emph{honki} — \textquotedblleft{}serious; in earnest\textquotedblright{}.

\begin{quote}
\textbf{\ja{ちちは} \ja{ほんきで} \ja{はなしたと} \ja{おもいます。}} — \emph{chichi wa honki de hanashita to omoimasu} — \textquotedblleft{}I think my father meant what he said.\textquotedblright{}
\end{quote}

A na-adjective: \ja{な} before a noun.

\begin{cousinweb}
\textbf{\ja{ほんき}} (\emph{honki}) is Sino-Japanese (kango): two Chinese characters read with their on'yomi, readings borrowed from Chinese. They are \emph{hon}, \textquotedblleft{}true\textquotedblright{} (Middle Chinese \emph{pwonX}) and \emph{ki}, \textquotedblleft{}spirit\textquotedblright{} (\emph{khj+jH}) (KANJIDIC2; Middle Chinese in Baxter's transcription).
\end{cousinweb}

\begin{grammarlens}[title={What you've built}]
One more word for giving reasons and opinions.
\end{grammarlens}

\subsection*{Your turn}
\begin{itemize}
\raggedright
  \item \emph{Say it:} \emph{honki}
  \item \emph{Say it:} \emph{honki}, once more
  \item \emph{Say it:} \emph{chichi wa honki de hanashita to omoimasu} — I think my father meant what he said.
  \item \emph{Recall:} say the Japanese for to admit, then the Japanese for to misunderstand, then say \emph{honki} again
\end{itemize}

\begin{tcolorbox}[breakable,colback=teal!4,colframe=teal!35!black,title={Before you move on}]
What does \ja{ほんき} mean? (\textquotedblleft{}Serious; in earnest\textquotedblright{}.) Say it once more.
\end{tcolorbox}
